\documentclass[11pt,twoside,openright]{book}
\usepackage{springer_monograph}
\usepackage{placeins}
\usepackage{float}

\hypersetup{
    pdftitle={Chapter 04: Reproducible Mass Spectrometry Project Architecture},
    pdfauthor={Lucas VHH TRAN}
}

\begin{document}

% --- Chapter Title Block ---
\setcounter{chapter}{3}
\chapter{Reproducible Mass Spectrometry Project Architecture}
\label{chap:reproducible_project}

\noindent
\textbf{Lucas VHH TRAN}\footnote{CK-CARE Davos, Christine K\"uhne--Center for Allergy Research and Education, Davos, Switzerland. Email: \texttt{lucas.tran@ck-care.ch}}\\
\rule{\textwidth}{0.8pt}

\begin{abstract}
\noindent
Reproducibility in mass spectrometry requires deterministic tracking of software versions, operating system dependencies, parameter configurations, and intermediate file caches \cite{peng2011reproducible, baker2016reproducibility, sandve2013ten}. Interactive graphical user interfaces obscure parameter histories and prevent exact replication across laboratories. In contrast, scripted pipelines define analytical operations as explicit code. A production-grade computational MS project relies on four architectural layers: portable file referencing that prevents hardcoded operating-system path failures \cite{muller2020here, marwick2018packaging}, dependency freezing via package lockfiles (\texttt{renv.lock}) \cite{ushey2024renv}, full-stack POSIX containerization (Docker) \cite{boettiger2015docker, nust2020ten}, and directed acyclic graph (DAG) caching (\texttt{targets}) \cite{landau2021targets}. DAG caching tracks cryptographic content hashes of raw spectral inputs and parameter values, recalculating only targets invalidated by code or parameter mutations while skipping computationally expensive peak-picking steps. This chapter presents the architecture of portable project layouts, cryptographic checksum verification, lockfile dependency auditing, DAG caching dynamics, and runtime contract assertions \cite{wickham2011testthat} in R 4.6.1.
\end{abstract}

\vspace{1em}
\noindent\textbf{Keywords:}\enspace Reproducibility, renv lockfile, targets DAG, Cryptographic hash, Pipeline caching, Docker containerization, Runtime assertions.

\vspace{1em}
\section*{Chapter Overview and Architectural Roadmap}

Table~\ref{tab:repro_roadmap} outlines the four tiers of the reproducibility stack for mass spectrometry research.

\begin{table}[htbp]
\centering
\caption{Four tiers of the reproducible mass spectrometry computing architecture.}
\label{tab:repro_roadmap}
\small
\begin{tabularx}{\textwidth}{>{\raggedright\arraybackslash}l >{\raggedright\arraybackslash}p{3.8cm} >{\raggedright\arraybackslash}X >{\raggedright\arraybackslash}p{2.8cm}}
\toprule
\textbf{Tier} & \textbf{Tool / Specification} & \textbf{Mechanism} & \textbf{Failure Mode Addressed} \\
\midrule
\textbf{1. Workspace} & RStudio Project / \texttt{here} \cite{muller2020here} & Anchor relative paths to root \texttt{.Rproj} directory & Hardcoded path breakage (\texttt{C:/Users/}) \\
\addlinespace
\textbf{2. Packages} & \texttt{renv} (\texttt{renv.lock}) \cite{ushey2024renv} & Freezes exact CRAN and Bioconductor commit versions & API deprecation and version drift \\
\addlinespace
\textbf{3. Operating System} & Docker / Rocker image \cite{boettiger2015docker, nust2020ten} & Bundles C/C++ system libraries (\texttt{libnetcdf}) & Dynamic link failures across Linux/OS \\
\addlinespace
\textbf{4. Pipeline DAG} & \texttt{targets} \cite{landau2021targets} & Cryptographic hash tracking of data and parameters & Redundant re-computation and stale figures \\
\bottomrule
\end{tabularx}
\end{table}

\begin{notebox}
\textbf{Tutor Warning: The "Works on My Machine" Fallacy}\\[0.4em]
An R script that executes successfully in an interactive R session can silently yield divergent results six months later if package versions or system shared libraries shift. An undocumented update to a peak detection package can alter peak boundary integration, changing quantitative matrices and downstream biological conclusions. Long-term reproducibility requires version-controlled lockfiles and automated pipeline graphs from the initial run.
\end{notebox}

\FloatBarrier
\section{Project Directory Topology and Path Portability}
\label{sec:directory_topology}

Computational mass spectrometry projects organize data files, transformation scripts, and reports into distinct directories \cite{marwick2018packaging, sandve2013ten}. Mixing raw vendor files with intermediate CSV outputs risks accidental modification of original instrument data.

\subsection{Standard Project Directory Layout}

A production mass spectrometry project uses a standardized directory layout anchored by an RStudio Project file:
\begin{verbatim}
ms_project_root/
  |-- ms_project.Rproj     # Anchors project root directory
  |-- data/
  |   |-- raw/             # Read-only vendor mzML / raw files
  |   \-- processed/       # SummarizedExperiment / RDS outputs
  |-- R/                   # Modular analysis functions
  |-- reports/             # Quarto (.qmd) documents and PDF/HTML output
  |-- _targets.R           # Pipeline DAG definition
  |-- _targets/            # Automated target cache (git-ignored)
  |-- renv.lock            # Exact R package dependency manifest
  |-- Dockerfile           # POSIX container definition
  \-- .gitignore           # Excludes raw data and local cache
\end{verbatim}

\subsection{Path Portability via Anchor Files}

Hardcoding absolute file paths (such as \texttt{"C:/Users/lucas/data/raw/run.mzML"}) breaks pipeline execution when a collaborator runs the code on a different computer or when the directory moves. The \texttt{here} package locates the project root dynamically by identifying the directory containing \texttt{.Rproj} or \texttt{.git} \cite{muller2020here}. It constructs platform-independent paths using forward slashes:
\begin{equation}
\text{path} = \operatorname{here}(\text{"data"}, \text{"raw"}, \text{"run\_01.mzML"})
\label{eq:here_path}
\end{equation}

Listing~4.1 computes MD5 and SHA-256 cryptographic hashes for an ingested raw spectral file.

\begin{rchunk}{Listing 4.1: Cryptographic Checksum Integrity Verification}
library(digest)

# Generate simulated raw mzML file in temporary workspace
temp_dir <- tempfile("ms_repro_audit_")
dir.create(temp_dir, recursive = TRUE)
raw_dir  <- file.path(temp_dir, "data", "raw")
dir.create(raw_dir, recursive = TRUE)

sim_raw_file <- file.path(raw_dir, "run_01.mzML")
writeLines(c("<?xml version='1.0' encoding='UTF-8'?>",
             "<mzML xmlns='http://psi.hupo.org/ms/mzml'>",
             "  <run id='sim_run_01' defaultInstrumentConfigurationRef='LCMS'>",
             "    <spectrum index='0' id='scan=1' defaultArrayLength='500'/>",
             "  </run>",
             "</mzML>"), con = sim_raw_file)

# Calculate MD5 and SHA-256 checksums
sha256_raw <- digest::digest(sim_raw_file, algo = "sha256", file = TRUE)
md5_raw    <- digest::digest(sim_raw_file, algo = "md5", file = TRUE)
bytes_raw  <- file.info(sim_raw_file)$size

cat(sprintf("Simulated Raw File: %s\n", basename(sim_raw_file)))
cat(sprintf("  File Size: %d bytes\n", bytes_raw))
cat(sprintf("  MD5 Checksum:    %s\n", md5_raw))
cat(sprintf("  SHA-256 Hash:    %s\n", sha256_raw))
\end{rchunk}

\begin{routput}
Simulated Raw File: run\_01.mzML\\
\enspace\enspace File Size: 233 bytes\\
\enspace\enspace MD5 Checksum:\enspace\enspace\enspace\enspace 10469785b424d205843712a94c037a76\\
\enspace\enspace SHA-256 Hash:\enspace\enspace\enspace\enspace 0e3f3fb807b1f1bad15147b3a8f356f4ba85962734d904a115ae3adfde352ace
\end{routput}

The raw \texttt{run\_01.mzML} file ($233$ bytes) produces SHA-256 hash \texttt{0e3f3fb8\dots}. If a byte changes during network transfer or disk storage, the cryptographic checksum changes immediately, exposing data corruption.

\FloatBarrier
\section{Software Environment Freezing: renv and Docker}
\label{sec:renv_docker}

R packages undergo continuous releases. Bioconductor coordinates releases every six months, while CRAN updates packages asynchronously. An analysis written against \texttt{xcms} version 3.20 may fail under version 4.0 if function arguments change \cite{beaulieujones2017reproducibility}.

\subsection{Package Version Pinning with renv}

The \texttt{renv} package establishes an isolated, project-local package library \cite{ushey2024renv}. The dependency graph is recorded in a JSON manifest (\texttt{renv.lock}). This file specifies the exact package versions, sources (CRAN, Bioconductor, or GitHub), and git commit SHAs:
\begin{equation}
\mathbf{E} = \{ (p_k, v_k, s_k, h_k) \}_{k=1}^K
\label{eq:renv_tuple}
\end{equation}
where $p_k$ is the package name, $v_k$ is the semantic version string, $s_k$ is the repository source, and $h_k$ is the commit hash.

Listing~4.2 parses and audits the live \texttt{renv.lock} file from the monograph repository.

\begin{rchunk}{Listing 4.2: renv Lockfile Parsing and Architecture Audit}
library(jsonlite)

proj_lock <- "renv.lock"
lock_data <- jsonlite::fromJSON(proj_lock)
r_version <- lock_data$R$Version
bioc_ver  <- lock_data$Bioconductor$Version
pkg_names <- names(lock_data$Packages)
n_pkgs    <- length(pkg_names)

repos <- vapply(lock_data$Packages, function(x) {
  if (!is.null(x$Repository)) x$Repository else "Other/GitHub"
}, character(1))
repo_counts <- table(repos)

cat(sprintf("Project renv.lock Found: %s\n", normalizePath(proj_lock)))
cat(sprintf("  Recorded R Version:            %s\n", r_version))
cat(sprintf("  Recorded Bioconductor Version: %s\n", bioc_ver))
cat(sprintf("  Total Locked Packages:         %d\n", n_pkgs))
cat("  Repository Distribution:\n")
for (repo in names(repo_counts)) {
  cat(sprintf("    - %s: %d packages (%.1f%%)\n", repo, repo_counts[repo],
              repo_counts[repo] / n_pkgs * 100))
}
\end{rchunk}

\begin{routput}
Project renv.lock Found: C:\textbackslash Users\textbackslash lucas\textbackslash OneDrive\textbackslash Statistics\textbackslash RforMS\textbackslash renv.lock\\
\enspace\enspace Recorded R Version:\enspace\enspace\enspace\enspace\enspace\enspace\enspace\enspace\enspace\enspace\enspace\enspace 4.6.1\\
\enspace\enspace Recorded Bioconductor Version: 3.23\\
\enspace\enspace Total Locked Packages:\enspace\enspace\enspace\enspace\enspace\enspace\enspace\enspace\enspace 159\\
\enspace\enspace Repository Distribution:\\
\enspace\enspace\enspace\enspace - Bioconductor 3.23: 9 packages (5.7\%)\\
\enspace\enspace\enspace\enspace - CRAN: 111 packages (69.8\%)\\
\enspace\enspace\enspace\enspace - https://bioc-release.r-universe.dev: 37 packages (23.3\%)\\
\enspace\enspace\enspace\enspace - https://bioc.r-universe.dev: 1 packages (0.6\%)\\
\enspace\enspace\enspace\enspace - Other/GitHub: 1 packages (0.6\%)
\end{routput}

The audit confirms 159 locked packages pinned to R 4.6.1 and Bioconductor 3.23. CRAN accounts for 111 packages ($69.8\%$), Bioconductor repositories account for 47 packages ($29.6\%$), and GitHub accounts for 1 package ($0.6\%$).

\subsection{Full-Stack POSIX Containerization (Docker)}

The \texttt{renv.lock} manifest pins R code, but R packages that parse mass spectrometry data depend on underlying C and C++ system libraries (such as NetCDF, HDF5, and XML2). Divergent system libraries can cause segmentation faults or numerical differences across operating systems.

Docker encapsulates the operating system, system libraries, R runtime, and package library into an immutable image \cite{boettiger2015docker, nust2020ten}. A standard container definition for mass spectrometry analysis builds from a Rocker base image:
\begin{verbatim}
FROM rocker/tidyverse:4.4.1
RUN apt-get update && apt-get install -y \
    libnetcdf-dev \
    libhdf5-dev \
    libxml2-dev
COPY . /workspace
WORKDIR /workspace
RUN R -e "renv::restore()"
CMD ["R", "-e", "targets::tar_make()"]
\end{verbatim}
Running this container reproduces the pipeline execution independently of host operating system configurations.

\FloatBarrier
\section{Directed Acyclic Graph (DAG) Pipeline Automation}
\label{sec:targets_dag}

Mass spectrometry pipelines involve computationally intensive stages: CentWave peak detection on 100 samples can require several hours. Re-running the entire pipeline after correcting a typographical error in a figure title is inefficient.

The \texttt{targets} package organizes an analysis into a Directed Acyclic Graph (DAG) \cite{landau2021targets}. Each step in the pipeline is defined as a target:
\begin{equation}
T_j = f_j(T_{j_1}, T_{j_2}, \dots, \theta_j)
\label{eq:dag_node}
\end{equation}
where $T_j$ depends on upstream targets and parameter values $\theta_j$.

\subsection{Cryptographic Target Invalidation}

The pipeline engine computes a cryptographic hash for each target:
\begin{equation}
H(T_j) = \operatorname{hash}\Big( H(T_{j_1}) \mathbin{\Vert} H(T_{j_2}) \mathbin{\Vert} \operatorname{hash}(f_j) \mathbin{\Vert} \operatorname{hash}(\theta_j) \Big)
\label{eq:target_hash}
\end{equation}
When \texttt{tar\_make()} executes:
\begin{enumerate}
\item If $H(T_j)$ matches the stored hash in the \texttt{\_targets/} cache, the step is skipped and the cached object is loaded from disk in $O(1)$ time.
\item If an upstream input or parameter $\theta_j$ changes, $H(T_j)$ changes, triggering re-computation of $T_j$ and all downstream dependent targets. Unrelated parallel branches remain untouched.
\end{enumerate}

Listing~4.3 simulates DAG hash generation, zero-cost cache skipping, and parameter invalidation.

\begin{rchunk}{Listing 4.3: DAG Pipeline Dependency Caching Simulation}
library(digest)

# Step 1: Compute initial DAG hashes for a 3-step MS pipeline
ppm_param <- 15.0
cache_store <- list()

hash_target1 <- digest::digest(list("raw_mzml_data"), algo = "xxhash64")
hash_target2 <- digest::digest(list(hash_target1, ppm = ppm_param), algo = "xxhash64")
hash_target3 <- digest::digest(list(hash_target2, method = "max_intensity"), algo = "xxhash64")

cat("Initial DAG Hash Tree Generation:\n")
cat(sprintf("  Target 1 [raw_file]:       %s\n", hash_target1))
cat(sprintf("  Target 2 [peak_detection]: %s (ppm = %.1f)\n", hash_target2, ppm_param))
cat(sprintf("  Target 3 [quant_matrix]:   %s\n", hash_target3))

# Populate cache
cache_store[[hash_target1]] <- list(status = "SKIPPED")
cache_store[[hash_target2]] <- list(status = "SKIPPED")
cache_store[[hash_target3]] <- list(status = "SKIPPED")

# Invocation 2: Zero input changes -> 100% cache hit
t_start <- Sys.time()
inv2_recompute <- c(
  t1 = is.null(cache_store[[hash_target1]]),
  t2 = is.null(cache_store[[hash_target2]]),
  t3 = is.null(cache_store[[hash_target3]])
)
t_inv2 <- as.numeric(difftime(Sys.time(), t_start, units = "secs"))

cat("\nInvocation 2 (Zero Input Changes):\n")
cat(sprintf("  Recomputed Targets: %d / 3 (Targets Skipped: %d)\n",
            sum(inv2_recompute), sum(!inv2_recompute)))
cat(sprintf("  Pipeline Verification Time: %.6f seconds (0 ms recompute penalty)\n", t_inv2))

# Invocation 3: Parameter mutation: ppm 15.0 -> 10.0
ppm_param_new <- 10.0
hash_target2_new <- digest::digest(list(hash_target1, ppm = ppm_param_new), algo = "xxhash64")
hash_target3_new <- digest::digest(list(hash_target2_new, method = "max_intensity"), algo = "xxhash64")

inv3_recompute <- c(
  t1_raw_file       = is.null(cache_store[[hash_target1]]),     # Cached!
  t2_peak_detection = is.null(cache_store[[hash_target2_new]]), # Invalidated!
  t3_quant_matrix   = is.null(cache_store[[hash_target3_new]])  # Downstream cascade!
)

cat("\nInvocation 3 (Parameter Mutation: ppm 15.0 -> 10.0):\n")
cat(sprintf("  Target 1 [raw_file]:       %s\n", ifelse(inv3_recompute["t1_raw_file"], "RECOMPUTED", "SKIPPED (Cached)")))
cat(sprintf("  Target 2 [peak_detection]: %s\n", ifelse(inv3_recompute["t2_peak_detection"], "RECOMPUTED (Invalidated)", "SKIPPED")))
cat(sprintf("  Target 3 [quant_matrix]:   %s\n", ifelse(inv3_recompute["t3_quant_matrix"], "RECOMPUTED (Downstream cascade)", "SKIPPED")))
cat(sprintf("  Total Recomputations: %d / 3 targets\n", sum(inv3_recompute)))
\end{rchunk}

\begin{routput}
Initial DAG Hash Tree Generation:\\
\enspace\enspace Target 1 [raw\_file]:\enspace\enspace\enspace\enspace\enspace\enspace 7ef71c205e4d1ac5\\
\enspace\enspace Target 2 [peak\_detection]: aa6b898bffea970e (ppm = 15.0)\\
\enspace\enspace Target 3 [quant\_matrix]:\enspace\enspace dd99504415ea3986\\[0.4em]
Invocation 2 (Zero Input Changes):\\
\enspace\enspace Recomputed Targets: 0 / 3 (Targets Skipped: 3)\\
\enspace\enspace Pipeline Verification Time: 0.000726 seconds (0 ms recompute penalty)\\[0.4em]
Invocation 3 (Parameter Mutation: ppm 15.0 -$>$ 10.0):\\
\enspace\enspace Target 1 [raw\_file]:\enspace\enspace\enspace\enspace\enspace\enspace SKIPPED (Cached)\\
\enspace\enspace Target 2 [peak\_detection]: RECOMPUTED (Invalidated)\\
\enspace\enspace Target 3 [quant\_matrix]:\enspace\enspace RECOMPUTED (Downstream cascade)\\
\enspace\enspace Total Recomputations: 2 / 3 targets
\end{routput}

In Invocation 2, verification completes in $0.000726$ seconds with zero recomputations. In Invocation 3, mutating the mass tolerance from $15.0$ to $10.0$ ppm invalidates Target 2 and cascades to Target 3, while preserving the raw file ingestion target (Target 1).

\FloatBarrier
\section{Defensive Runtime Contract Testing on MS Assay Matrices}
\label{sec:contract_testing}

Silent data corruption occurs when erroneous preprocessing produces negative intensities, NaN values, or unexpected dimension drops without halting execution. Runtime contract assertions using \texttt{testthat} \cite{wickham2011testthat} verify mathematical invariants before data enter statistical models.

Listing~4.4 tests three structural contracts on a simulated feature abundance matrix: dimension consistency, intensity non-negativity, and missingness rate bounds.

\begin{rchunk}{Listing 4.4: Unit Testing and Defensive Contract Auditing}
library(testthat)

set.seed(99)
test_assay <- matrix(abs(rnorm(100, mean = 5000, sd = 1000)), nrow = 20, ncol = 5)
colnames(test_assay) <- paste0("Sample_", 1:5)
rownames(test_assay) <- paste0("Feat_", 1:20)
test_assay[1, 2] <- NA_real_

# Execute contract assertions
test_results <- testthat::test_that("MS Assay Integrity Contracts", {
  # Contract 1: Dimensional invariance
  testthat::expect_equal(dim(test_assay), c(20L, 5L))
  # Contract 2: Intensity non-negativity
  testthat::expect_true(all(test_assay[!is.na(test_assay)] >= 0))
  # Contract 3: Missingness upper bound (< 10%)
  missing_pct <- mean(is.na(test_assay)) * 100
  testthat::expect_lt(missing_pct, 10.0)
})

cat("Testthat Contract Suite Execution: 3 / 3 assertions passed.\n")
cat(sprintf("  Matrix Dimensions: %d features x %d samples\n", nrow(test_assay), ncol(test_assay)))
cat(sprintf("  Minimum Valid Intensity: %.2f counts\n", min(test_assay, na.rm = TRUE)))
cat(sprintf("  Observed Missingness Rate: %.2f%% (Contract threshold: < 10.00%%)\n",
            mean(is.na(test_assay)) * 100))
\end{rchunk}

\begin{routput}
Testthat Contract Suite Execution: 3 / 3 assertions passed.\\
\enspace\enspace Matrix Dimensions: 20 features x 5 samples\\
\enspace\enspace Minimum Valid Intensity: 1959.07 counts\\
\enspace\enspace Observed Missingness Rate: 1.00\% (Contract threshold: $<$ 10.00\%)
\end{routput}

All three contracts pass: the matrix contains $20$ features and $5$ samples, the minimum observed intensity is $1{,}959.07$ counts, and the observed missingness rate is $1.00\%$, satisfying the $< 10.00\%$ contract threshold.

\FloatBarrier
\section{Worked Exercises}
\label{sec:worked_exercises}

\subsection*{Exercise 4.1: Portable Path Resolution with \texttt{here::here()}}
\noindent\textbf{Problem Statement:}\enspace
Explain why hardcoding paths such as \texttt{"C:/Users/name/data/raw/sample.mzML"} causes pipeline failures across collaborators' systems and explain how \texttt{here::here()} eliminates this failure mode.

\noindent\textbf{Solution:}\enspace
Hardcoded paths contain user-specific root directories (\texttt{"C:/Users/name/"}) and Windows backslash separators that do not exist on other machines or POSIX operating systems (Linux/macOS). When another user runs the script, the file cannot be located, halting execution.

The function \texttt{here::here("data", "raw", "sample.mzML")} dynamically detects the project root by searching upward from the current working directory for an anchor file (such as \texttt{.Rproj} or \texttt{.git}). It joins directory components using the host operating system's native file separator. Moving the project directory to another disk, computer, or cloud container preserves valid path resolution without altering code.

\subsection*{Exercise 4.2: Dependency Tier Scope: \texttt{renv} vs Docker}
\noindent\textbf{Problem Statement:}\enspace
Distinguish the scope of reproducibility provided by \texttt{renv} from that provided by Docker. State a failure scenario where \texttt{renv.lock} is insufficient to guarantee identical numerical outputs.

\noindent\textbf{Solution:}\enspace
\begin{enumerate}
\item \textbf{\texttt{renv}}: Freezes the R package layer. It records exact package version strings, repository origins, and GitHub commit SHAs in \texttt{renv.lock}, ensuring that identical R bytecode is loaded.
\item \textbf{Docker}: Freezes the complete POSIX operating system layer. It bundles the Linux distribution kernel interface, compilers, system libraries (\texttt{libnetcdf}, \texttt{libxml2}, \texttt{libhdf5}), and BLAS/LAPACK linear algebra implementations.
\end{enumerate}

\noindent\textbf{Failure Scenario:}\enspace
An analyst runs \texttt{xcms} on Linux with NetCDF version 4.8, while a collaborator restores the identical \texttt{renv.lock} on macOS with NetCDF version 4.9. If NetCDF 4.9 introduces a slight difference in floating-point decompression of vendor raw arrays, downstream peak integration boundaries will diverge, producing altered quantitative values despite identical R package versions. Docker eliminates this divergence by freezing the compiled C libraries.

\subsection*{Exercise 4.3: DAG Invalidation Dynamics in \texttt{targets}}
\noindent\textbf{Problem Statement:}\enspace
In a mass spectrometry \texttt{targets} pipeline comprising:
\begin{enumerate}
\item \texttt{raw\_files}: Reads list of 50 mzML files.
\item \texttt{ppm\_tolerance}: Target defined as numeric scalar $15.0$.
\item \texttt{chrom\_peaks}: Calls \texttt{findChromPeaks(raw\_files, ppm = ppm\_tolerance)}.
\item \texttt{quant\_matrix}: Extracts feature intensity values from \texttt{chrom\_peaks}.
\item \texttt{report}: Renders Quarto report displaying the volcano plot.
\end{enumerate}
If the analyst alters \texttt{ppm\_tolerance} from $15.0$ to $10.0$, identify which targets \texttt{tar\_make()} skips and which targets it recomputes.

\noindent\textbf{Solution:}\enspace
\begin{itemize}
\item \textbf{\texttt{raw\_files}}: \textbf{SKIPPED (Cached)}. The raw mzML files and their file-level checksums are unmodified; their cryptographic hash is identical.
\item \textbf{\texttt{ppm\_tolerance}}: \textbf{RECOMPUTED}. The target definition mutated from $15.0$ to $10.0$.
\item \textbf{\texttt{chrom\_peaks}}: \textbf{RECOMPUTED}. This target depends on \texttt{ppm\_tolerance}; the change in upstream hash invalidates the cached peak table.
\item \textbf{\texttt{quant\_matrix}}: \textbf{RECOMPUTED}. Depends on the newly recomputed \texttt{chrom\_peaks}.
\item \textbf{\texttt{report}}: \textbf{RECOMPUTED}. Depends on \texttt{quant\_matrix}; the Quarto report is re-rendered with updated quantification results.
\end{itemize}

\subsection*{Exercise 4.4: High-Performance Serialization Formats}
\noindent\textbf{Problem Statement:}\enspace
Explain why default R serialization via \texttt{saveRDS()} creates bottlenecks in large mass spectrometry pipelines and state two alternative serialization formats supported by \texttt{targets}.

\noindent\textbf{Solution:}\enspace
Default \texttt{saveRDS()} uses single-threaded gzip compression. For large mass spectrometry objects ($> 2$ GB), gzip compression and decompression require substantial CPU time and inflate RAM during serialization.

Two high-performance alternatives are:
\begin{enumerate}
\item \textbf{\texttt{qs} (Quick Serialization)}: Uses multi-threaded Zstandard (zstd) compression, providing up to $5\times$ faster read and write throughput than \texttt{saveRDS()} with comparable compression ratios.
\item \textbf{\texttt{fst} (Fast Serialization of DataFrames)}: Provides multi-threaded, memory-mapped random-access storage for tabular feature matrices and PSM tables, enabling instant decompression of specific column subsets without loading the entire matrix into RAM.
\end{enumerate}

\FloatBarrier
\begin{thebibliography}{18}

\bibitem{landau2021targets}
Landau, W.M.
\newblock The targets {R} package: a dynamic Make-like function-oriented pipeline toolkit for reproducible workflows.
\newblock \emph{Journal of Open Source Software}, 6(57):2959, 2021.
\newblock \url{https://doi.org/10.21105/joss.02959}.

\bibitem{ushey2024renv}
Ushey, K., Wickham, H.
\newblock \emph{renv: Project Environments for R}.
\newblock R package version 1.0.7, 2024.
\newblock \url{https://CRAN.R-project.org/package=renv}.

\bibitem{peng2011reproducible}
Peng, R.D.
\newblock Reproducible research in computational science.
\newblock \emph{Science}, 334(6060):1226--1227, 2011.
\newblock \url{https://doi.org/10.1126/science.1213847}.

\bibitem{baker2016reproducibility}
Baker, M.
\newblock 1,500 scientists lift the lid on reproducibility.
\newblock \emph{Nature}, 533(7604):452--454, 2016.
\newblock \url{https://doi.org/10.1038/533452a}.

\bibitem{sandve2013ten}
Sandve, G.K., Nekrutenko, A., Taylor, J., Hovig, E.
\newblock Ten simple rules for reproducible computational research.
\newblock \emph{PLoS Computational Biology}, 9(10):e1003285, 2013.
\newblock \url{https://doi.org/10.1371/journal.pcbi.1003285}.

\bibitem{marwick2018packaging}
Marwick, B., Boettiger, C., Mullen, L.
\newblock Packaging data analytical work reproducibly using {R} packages, renv, and {Docker}.
\newblock \emph{The American Statistician}, 72(1):80--88, 2018.
\newblock \url{https://doi.org/10.1080/00031305.2017.1375986}.

\bibitem{muller2020here}
M\"uller, K.
\newblock \emph{here: A Simpler Way to Find Your Files}.
\newblock R package version 1.0.1, 2020.
\newblock \url{https://CRAN.R-project.org/package=here}.

\bibitem{boettiger2015docker}
Boettiger, C.
\newblock An introduction to {Docker} for reproducible research.
\newblock \emph{ACM SIGOPS Operating Systems Review}, 49(1):71--79, 2015.
\newblock \url{https://doi.org/10.1145/2723872.2723882}.

\bibitem{nust2020ten}
N\"ust, D., Sochat, V., Marwick, B., Eglen, S.J., Head, T., Hirst, T., Evans, B.D.
\newblock Ten simple rules for writing {Dockerfiles} for reproducible data science.
\newblock \emph{PLoS Computational Biology}, 16(11):e1008316, 2020.
\newblock \url{https://doi.org/10.1371/journal.pcbi.1008316}.

\bibitem{beaulieujones2017reproducibility}
Beaulieu-Jones, B.K., Greene, C.S.
\newblock Reproducibility of computational workflows is automated using continuous analysis.
\newblock \emph{Nature Biotechnology}, 35(4):342--346, 2017.
\newblock \url{https://doi.org/10.1038/nbt.3780}.

\bibitem{knuth1984literate}
Knuth, D.E.
\newblock Literate programming.
\newblock \emph{The Computer Journal}, 27(2):97--111, 1984.
\newblock \url{https://doi.org/10.1093/comjnl/27.2.97}.

\bibitem{wickham2011testthat}
Wickham, H.
\newblock testthat: Get started with testing.
\newblock \emph{The R Journal}, 3(1):5--10, 2011.
\newblock \url{https://doi.org/10.32614/RJ-2011-002}.

\bibitem{wilkinson2016fair}
Wilkinson, M.D., Dumontier, M., Aalbersberg, I.J., et~al.
\newblock The {FAIR} Guiding Principles for scientific data management and stewardship.
\newblock \emph{Scientific Data}, 3:160018, 2016.
\newblock \url{https://doi.org/10.1038/sdata.2016.18}.

\bibitem{deutsch2020pride}
Deutsch, E.W., Bandeira, N., Sharma, V., et~al.
\newblock The {ProteomeXchange} consortium in 2020: enabling 'big data' approaches in proteomics.
\newblock \emph{Nucleic Acids Research}, 48(D1):D1145--D1152, 2020.
\newblock \url{https://doi.org/10.1093/nar/gkz984}.

\bibitem{perezriverol2022pride}
Perez-Riverol, Y., Bai, J., Bandla, C., et~al.
\newblock The {PRIDE} database resources in 2022: a hub for mass spectrometry-based proteomics evidences.
\newblock \emph{Nucleic Acids Research}, 50(D1):D543--D552, 2022.
\newblock \url{https://doi.org/10.1093/nar/gkab1038}.

\bibitem{gatto2021msnbase}
Gatto, L., Gibb, S., Rainer, J.
\newblock {MSnbase}, efficient and elegant {R}-based processing and visualization of raw mass spectrometry data.
\newblock \emph{Journal of Proteome Research}, 20(1):1063--1069, 2021.
\newblock \url{https://doi.org/10.1021/acs.jproteome.0c00313}.

\bibitem{rainer2022spectra}
Rainer, J., Vicini, A., Salzer, L., et~al.
\newblock A modular and expandable ecosystem for metabolomics data annotation in {R}.
\newblock \emph{Metabolites}, 12(2):173, 2022.
\newblock \url{https://doi.org/10.3390/metabo12020173}.

\bibitem{blankenberg2010galaxy}
Blankenberg, D., Kuster, G.V., Coraor, N., et~al.
\newblock Galaxy: a web-based genome analysis tool for experimentalists.
\newblock \emph{Current Protocols in Molecular Biology}, 89(1):19--10, 2010.
\newblock \url{https://doi.org/10.1002/0471142727.mb1910s89}.

\end{thebibliography}

\end{document}
